\documentclass[12pt, letterpaper]{article}

% ── Paquetes ──────────────────────────────────────────────────────────────────
\usepackage[utf8]{inputenc}
\usepackage[T1]{fontenc}
\usepackage[spanish]{babel}
\usepackage{geometry}
\usepackage{booktabs}
\usepackage{array}
\usepackage{longtable}
\usepackage{xcolor}
\usepackage{colortbl}
\usepackage{fancyhdr}
\usepackage{titlesec}
\usepackage{enumitem}
\usepackage{hyperref}
\usepackage{graphicx}
\usepackage{tabularx}
\usepackage{multirow}
\usepackage{amsmath}
\usepackage{lmodern}
\usepackage{microtype}

% ── Márgenes ──────────────────────────────────────────────────────────────────
\geometry{
  top        = 2.5cm,
  bottom     = 2.5cm,
  left       = 3cm,
  right      = 2.5cm,
  headheight = 14pt
}

% ── Colores ───────────────────────────────────────────────────────────────────
\definecolor{azulSAR}{RGB}{29, 78, 216}
\definecolor{grisCabecera}{RGB}{243, 244, 246}
\definecolor{grisLinea}{RGB}{209, 213, 219}
\definecolor{azulClaro}{RGB}{219, 234, 254}

% ── Encabezado y pie de página ────────────────────────────────────────────────
\pagestyle{fancy}
\fancyhf{}
\fancyhead[L]{\small\textcolor{azulSAR}{\textbf{DEVS MRS S. DE R. L.}}}
\fancyhead[R]{\small\textcolor{gray}{Especificaciones Técnicas --- Sistema de Facturación}}
\fancyfoot[C]{\small\thepage}
\fancyfoot[R]{\small\textcolor{gray}{Uso interno / SAR Honduras}}
\renewcommand{\headrulewidth}{0.4pt}
\renewcommand{\footrulewidth}{0.4pt}

% ── Títulos de sección ────────────────────────────────────────────────────────
\titleformat{\section}{\large\bfseries\color{azulSAR}}{}{0em}{}[\titlerule]
\titleformat{\subsection}{\normalsize\bfseries\color{azulSAR}}{}{0em}{}

% ── Hipervínculos ─────────────────────────────────────────────────────────────
\hypersetup{
  colorlinks = true,
  linkcolor  = azulSAR,
  urlcolor   = azulSAR,
}

% ─────────────────────────────────────────────────────────────────────────────
\begin{document}

% ── PORTADA ───────────────────────────────────────────────────────────────────
\begin{titlepage}
  \centering
  \vspace*{2cm}

  {\Huge\bfseries\color{azulSAR} DEVS MRS S. DE R. L.\par}
  \vspace{0.5cm}
  {\large Col.\ Las Torres, Contiguo a la Puma B34 S1\\
         Tegucigalpa, M.D.C., Honduras\par}

  \vspace{2cm}
  \rule{0.8\linewidth}{1.5pt}
  \vspace{0.4cm}

  {\LARGE\bfseries Especificaciones Técnicas\\[0.3cm]
   Sistema Electrónico de Facturación\par}

  \vspace{0.4cm}
  \rule{0.8\linewidth}{1.5pt}

  \vspace{2cm}
  {\large\textbf{Documento preparado para:}\\[0.3cm]
   \textbf{Servicio de Administración de Rentas (SAR)}\\
   República de Honduras\par}

  \vfill

  \begin{tabular}{ll}
    \textbf{Versión:}       & 2.0                    \\[4pt]
    \textbf{Fecha:}         & Mayo 2026              \\[4pt]
    \textbf{Clasificación:} & Uso Interno / Registro \\
  \end{tabular}
\end{titlepage}

% ── TABLA DE CONTENIDOS ───────────────────────────────────────────────────────
\tableofcontents
\newpage

% ─────────────────────────────────────────────────────────────────────────────
\section{Información General de la Empresa}

\begin{tabularx}{\linewidth}{>{\bfseries}lX}
  \toprule
  Razón Social          & DEVS MRS S. DE R. L.                              \\
  Dirección             & Col.\ Las Torres, Contiguo a la Puma B34 S1,
                          Tegucigalpa, M.D.C., Honduras                      \\
  RTN                   & (Número de RTN de la empresa)                     \\
  Actividad Económica   & Servicios de Desarrollo de Software y Tecnología  \\
  Régimen Tributario    & Régimen General                                   \\
  Obligación Tributaria & ISV 15\% --- Impuesto Sobre Ventas                \\
  Moneda                & Lempira Hondureño (HNL)                           \\
  \bottomrule
\end{tabularx}

% ─────────────────────────────────────────────────────────────────────────────
\section{Descripción del Sistema}

El sistema es una aplicación web integral de gestión comercial y facturación
desarrollada internamente por DEVS MRS S.\ de R.\ L., diseñada para cumplir
con los requerimientos del formato de factura autorizado por el SAR de Honduras.
Permite la emisión de \textbf{Facturas Fiscales} y \textbf{Cotizaciones},
con flujo completo de gestión de clientes, servicios, gastos y reportería
financiera.

\subsection{Módulos del Sistema}

\begin{enumerate}[label=\textbf{\arabic*.}]
  \item \textbf{Facturación} --- Emisión de facturas fiscales con número
        correlativo, CAI, RTN, rangos autorizados e ISV desglosado.
  \item \textbf{Cotizaciones} --- Documento pre-venta con flujo de aprobación
        por correo electrónico; puede convertirse automáticamente en Factura
        al ser aprobada por el cliente.
  \item \textbf{Catálogo de Clientes} --- Registro centralizado de clientes
        con nombre, razón social, RTN, dirección, correo y teléfono.
        Vinculado al ecosistema de cotizaciones y facturas.
  \item \textbf{Catálogo de Servicios} --- Inventario de servicios tecnológicos
        ofrecidos, organizados por categorías, con precios y descripciones.
  \item \textbf{Gastos} --- Registro de gastos operativos categorizados para
        control financiero interno.
  \item \textbf{Reportería} --- Informes de ventas por fecha, cliente y
        servicio; resumen financiero del período.
  \item \textbf{Finanzas} --- Resumen financiero consolidado (ingresos,
        gastos, utilidad bruta).
  \item \textbf{Configuración} --- Parámetros de empresa, datos fiscales
        (CAI, rangos, RTN) y configuración de correo saliente (SMTP).
  \item \textbf{Gestión de Usuarios} --- Control de acceso con múltiples
        usuarios y autenticación JWT.
\end{enumerate}

% ─────────────────────────────────────────────────────────────────────────────
\section{Arquitectura Tecnológica}

\subsection{Tecnologías Utilizadas}

\begin{tabularx}{\linewidth}{>{\bfseries}l l X}
  \toprule
  \rowcolor{azulClaro}
  Capa              & Tecnología              & Versión / Descripción        \\
  \midrule
  Backend           & Node.js                 & v22 LTS (runtime JavaScript) \\
  Framework API     & Express.js              & v4 (API REST HTTP)           \\
  Base de Datos     & NeDB (nedb-promises)    & Base de datos NoSQL embebida,
                                               archivos \texttt{.db} en disco \\
  Autenticación     & JSON Web Tokens (JWT)   & Tokens firmados para sesión  \\
  Email             & Nodemailer              & SMTP multi-proveedor         \\
  Email OAuth2      & Microsoft Graph API     & Outlook personal vía OAuth2  \\
  Frontend          & React                   & v18                          \\
  Bundler           & Vite                    & v5                           \\
  Estilos           & Tailwind CSS            & v3                           \\
  Lenguaje          & JavaScript (ES Modules) & ECMAScript 2022+             \\
  Sistema Operativo & Windows / Linux         & Multiplataforma              \\
  \bottomrule
\end{tabularx}

\subsection{Arquitectura de la Aplicación}

\begin{itemize}
  \item \textbf{Tipo:} Aplicación web cliente-servidor (SPA + REST API).
  \item \textbf{Comunicación:} HTTP/1.1, JSON en todos los endpoints.
  \item \textbf{Autenticación:} JWT Bearer Token en cabecera \texttt{Authorization}.
  \item \textbf{Puerto API:} 3001 (configurable mediante variable de entorno).
  \item \textbf{Puerto Frontend:} 5173 en desarrollo; se puede servir de
        forma estática en producción.
  \item \textbf{Despliegue:} Servidor local o nube privada; no requiere
        base de datos externa.
  \item \textbf{Rutas públicas:} El endpoint de respuesta a cotizaciones
        (\texttt{/api/cotizacion-respuesta/:token}) no requiere autenticación,
        permitiendo al cliente aprobar o rechazar desde el correo recibido.
\end{itemize}

% ─────────────────────────────────────────────────────────────────────────────
\section{Base de Datos}

\subsection{Motor de Almacenamiento}

Se utiliza \textbf{NeDB} (Node Embedded Database), una base de datos NoSQL
embebida que almacena los datos en archivos de texto plano en disco
(\texttt{.db}). No requiere instalación de servidor de base de datos externo.

\subsection{Colecciones (Datastores)}

\begin{tabularx}{\linewidth}{>{\bfseries}l X}
  \toprule
  Colección              & Descripción                                        \\
  \midrule
  \texttt{invoices.db}   & Facturas y cotizaciones con todos sus campos       \\
  \texttt{clientes.db}   & Catálogo de clientes (nombre, empresa, RTN, etc.)  \\
  \texttt{servicios.db}  & Catálogo de servicios tecnológicos ofrecidos       \\
  \texttt{gastos.db}     & Registro de gastos operativos                      \\
  \texttt{users.db}      & Usuarios del sistema con contraseñas cifradas      \\
  \texttt{config.db}     & Configuración de empresa (CAI, RTN, SMTP, etc.)    \\
  \bottomrule
\end{tabularx}

\subsection{Ruta de Almacenamiento}

\begin{verbatim}
backend/data/invoices.db
backend/data/clientes.db
backend/data/servicios.db
backend/data/gastos.db
backend/data/users.db
backend/data/config.db
\end{verbatim}

% ─────────────────────────────────────────────────────────────────────────────
\section{Formato de Numeración de Documentos}

\subsection{Factura Fiscal}

El número de factura sigue el formato oficial establecido por el SAR:

\begin{center}
  \Large\texttt{000-001-01-XXXXXXXX}
\end{center}

\begin{tabularx}{\linewidth}{>{\ttfamily\bfseries}l X}
  \toprule
  Segmento   & Significado                                       \\
  \midrule
  000        & Identificador de establecimiento                  \\
  001        & Punto de emisión                                  \\
  01         & Tipo de documento (01 = Factura)                  \\
  XXXXXXXX   & Número correlativo de 8 dígitos (00000001$\ldots$)\\
  \bottomrule
\end{tabularx}

\subsection{Cotización}

\begin{center}
  \Large\texttt{COT-XXXXXX}
\end{center}

Las cotizaciones usan el prefijo \texttt{COT-} seguido de un número
correlativo de 6 dígitos, independiente del correlativo de facturas.

% ─────────────────────────────────────────────────────────────────────────────
\section{Campos del Documento de Factura / Cotización}

\subsection{Datos del Emisor (Empresa)}

\begin{tabularx}{\linewidth}{>{\bfseries}l X}
  \toprule
  Campo                     & Descripción                              \\
  \midrule
  Razón Social              & Nombre de la empresa emisora             \\
  Dirección                 & Dirección fiscal de la empresa           \\
  Teléfono                  & Número de teléfono                       \\
  Correo Electrónico        & Email de contacto                        \\
  RTN                       & Registro Tributario Nacional de la empresa\\
  CAI                       & Código de Autorización de Impresión
                              (emitido por el SAR) --- solo en facturas  \\
  Rango Autorizado (Desde)  & Primer número de factura autorizado      \\
  Rango Autorizado (Hasta)  & Último número de factura autorizado      \\
  Fecha Límite de Emisión   & Fecha máxima para emitir dentro del rango\\
  \bottomrule
\end{tabularx}

\subsection{Datos del Cliente (Adquirente)}

\begin{tabularx}{\linewidth}{>{\bfseries}l X}
  \toprule
  Campo          & Descripción                                          \\
  \midrule
  Señor(es)      & Razón social de la empresa cliente; si no hay empresa,
                   se muestra el nombre del contacto                     \\
  Razón Social   & Nombre de la empresa u organización del cliente      \\
  Dirección      & Dirección del cliente                                \\
  RTN            & Registro Tributario Nacional del cliente (opcional)  \\
  Correo         & Correo electrónico del cliente (para envío de
                   cotizaciones y notificaciones)                        \\
  \bottomrule
\end{tabularx}

\subsection{Datos Fiscales del Documento}

\begin{tabularx}{\linewidth}{>{\bfseries}l X}
  \toprule
  Campo                         & Descripción                          \\
  \midrule
  Número de Documento           & Correlativo en formato SAR (factura)
                                  o COT-XXXXXX (cotización)             \\
  Fecha                         & Fecha de emisión                     \\
  Lugar de Emisión              & Tegucigalpa, M.D.C.                  \\
  No.\ O/C Exenta               & Orden de compra exenta (si aplica)   \\
  No.\ Registro de Exonerado    & Registro de exoneración              \\
  No.\ Registro de la SAG       & Registro SAG (sector agropecuario)   \\
  \bottomrule
\end{tabularx}

\subsection{Detalle de Ítems}

\begin{tabularx}{\linewidth}{>{\bfseries}l X}
  \toprule
  Campo          & Descripción                                         \\
  \midrule
  Cantidad       & Número de unidades                                  \\
  Descripción    & Descripción del bien o servicio                     \\
  Precio Unitario& Precio por unidad en Lempiras                       \\
  Valor Total    & Cantidad $\times$ Precio Unitario                   \\
  \bottomrule
\end{tabularx}

\subsection{Totales y Cálculos Fiscales}

\begin{tabularx}{\linewidth}{>{\bfseries}l X}
  \toprule
  Campo                        & Descripción                            \\
  \midrule
  Venta Exonerada              & Monto exento de ISV por exoneración    \\
  Sub-Total Exento             & Subtotal de bienes/servicios exentos   \\
  Descuentos y Rebajas         & Total de descuentos aplicados          \\
  Sub-Total Gravado            & Base imponible para ISV                \\
  15\% Impuesto Sobre Ventas   & ISV calculado sobre el gravado         \\
  \textbf{TOTAL}               & Total a pagar en Lempiras              \\
  La Suma de                   & Total en letras (Lempiras)             \\
  \bottomrule
\end{tabularx}

% ─────────────────────────────────────────────────────────────────────────────
\section{Cálculo del ISV (Impuesto Sobre Ventas)}

El sistema soporta dos modalidades de cálculo del ISV del 15\%:

\subsection{Modalidad 1: ISV No Incluido en el Precio}

El precio ingresado \textbf{no} incluye el impuesto; el ISV se agrega al total.

\begin{align}
  \text{Suma Bruta}        &= \sum_{i} (\text{Cantidad}_i \times \text{PrecioUnit}_i) \\
  \text{SubTotalGravado}   &= \text{SumaBruta} - \text{VentaExonerada}
                              - \text{SubtotalExento} - \text{Descuentos}             \\
  \text{ISV}               &= \text{SubTotalGravado} \times 0.15                      \\
  \text{TOTAL}             &= \text{VentaExonerada} + \text{SubtotalExento}
                              + \text{SubTotalGravado} + \text{ISV}
\end{align}

\subsection{Modalidad 2: ISV Incluido en el Precio}

El precio ingresado \textbf{ya incluye} el ISV; se extrae la base gravable.

\begin{align}
  \text{BrutoConISV}       &= \text{SumaBruta} - \text{VentaExonerada}
                              - \text{SubtotalExento} - \text{Descuentos}             \\
  \text{SubTotalGravado}   &= \frac{\text{BrutoConISV}}{1.15}                         \\
  \text{ISV}               &= \text{BrutoConISV} - \text{SubTotalGravado}             \\
  \text{TOTAL}             &= \text{VentaExonerada} + \text{SubtotalExento}
                              + \text{BrutoConISV}
\end{align}

\textit{Nota: El factor divisor $1.15$ equivale a $1 + 0.15$, cumpliendo
la norma de desglose de ISV incluido.}

% ─────────────────────────────────────────────────────────────────────────────
\section{Flujo de Documentos}

\subsection{Factura Directa}

\begin{center}
  \textbf{Nueva Factura} $\longrightarrow$ \textbf{Emitida}
  $\longrightarrow$ \textit{[Anulada (opcional)]}
\end{center}

\subsection{Flujo Cotización \texorpdfstring{$\rightarrow$}{→} Factura}

\begin{center}
  \textbf{Nueva Cotización (pendiente)}
  $\longrightarrow$ \textbf{Enviada por correo al cliente}
  $\longrightarrow$ \textbf{Aprobada / Rechazada por el cliente}
  $\longrightarrow$ \textbf{Factura generada automáticamente}
\end{center}

\begin{itemize}
  \item La cotización queda en estado \textit{facturada} y almacena la
        referencia al número de factura generado.
  \item La factura resultante almacena \texttt{from\_cotizacion\_number}
        para trazabilidad.
  \item No se puede editar una cotización ya \textit{facturada}.
  \item El cliente también puede aprobar o rechazar manualmente desde
        el sistema interno sin necesidad de envío por correo.
\end{itemize}

% ─────────────────────────────────────────────────────────────────────────────
\section{Flujo de Aprobación por Correo Electrónico}

\subsection{Descripción del Flujo}

\begin{enumerate}[label=\textbf{\arabic*.}]
  \item El usuario envía la cotización al cliente por correo desde el
        sistema, especificando el correo del destinatario.
  \item El sistema genera un \textbf{token de aprobación único} (validez
        de 7 días) y lo almacena junto con la cotización.
  \item Se envía un correo HTML al cliente con:
    \begin{itemize}
      \item Nombre/razón social del cliente y RTN.
      \item Número de cotización, fecha y total.
      \item Tabla de servicios cotizados con precios.
      \item Enlace \textbf{``Ver Cotización Completa / Imprimir PDF''}.
      \item Botones de \textbf{Aprobar} y \textbf{Rechazar}.
    \end{itemize}
  \item El cliente abre el enlace y visualiza el documento formal completo
        con el diseño oficial de la empresa.
  \item Al aprobar, el sistema genera automáticamente la factura y notifica
        al usuario interno.
  \item Al rechazar, el cliente ingresa el motivo; el sistema registra la
        razón de rechazo en la cotización.
\end{enumerate}

\subsection{Proveedores de Correo Soportados}

\begin{tabularx}{\linewidth}{>{\bfseries}l X}
  \toprule
  Proveedor          & Método de Autenticación                             \\
  \midrule
  Gmail              & App Password (contraseña de aplicación de Google)   \\
  Yahoo              & Contraseña / App Password                           \\
  SMTP Personalizado & Host, puerto, usuario y contraseña configurables    \\
  Outlook Personal   & OAuth2 via Microsoft Graph API (Device Code Flow)   \\
  \bottomrule
\end{tabularx}

\subsection{Nota sobre Outlook Personal}

Las cuentas personales de Outlook.com (\texttt{@outlook.com},
\texttt{@hotmail.com}) no soportan autenticación SMTP básica desde 2023.
El sistema implementa el flujo \textbf{OAuth2 Device Code} mediante
Microsoft Graph API, que requiere el registro de una aplicación gratuita
en Azure Portal (\texttt{portal.azure.com}) con soporte para cuentas
personales de Microsoft.

% ─────────────────────────────────────────────────────────────────────────────
\section{API REST del Sistema}

\begin{longtable}{>{\ttfamily\small\bfseries}l l X}
  \toprule
  \normalfont\textbf{Endpoint}          & \textbf{Método} & \textbf{Descripción} \\
  \midrule
  \endhead
  \multicolumn{3}{l}{\textit{Autenticación (públicas)}} \\
  /api/auth/login                       & POST   & Iniciar sesión, retorna JWT        \\
  /api/auth/users                       & GET    & Listar usuarios del sistema        \\
  /api/auth/users                       & POST   & Crear nuevo usuario                \\
  /api/auth/users/:id                   & DELETE & Eliminar usuario                   \\
  /api/auth/users/:id/password          & PUT    & Cambiar contraseña                 \\
  \midrule
  \multicolumn{3}{l}{\textit{Facturas y Cotizaciones (protegidas)}} \\
  /api/invoices                         & GET    & Listar documentos (filtros: tipo, fecha) \\
  /api/invoices                         & POST   & Crear factura o cotización         \\
  /api/invoices/:id                     & GET    & Obtener documento por ID           \\
  /api/invoices/:id                     & PUT    & Editar cotización pendiente        \\
  /api/invoices/:id/status              & PATCH  & Cambiar estado de cotización       \\
  /api/invoices/:id/convert             & POST   & Convertir cotización a factura     \\
  /api/invoices/:id                     & DELETE & Anular documento                   \\
  /api/invoices/:id/enviar-cotizacion   & POST   & Enviar cotización por correo       \\
  \midrule
  \multicolumn{3}{l}{\textit{Clientes (protegidas)}} \\
  /api/clientes                         & GET    & Listar clientes                    \\
  /api/clientes                         & POST   & Crear cliente                      \\
  /api/clientes/:id                     & GET    & Obtener cliente                    \\
  /api/clientes/:id                     & PUT    & Actualizar cliente                 \\
  /api/clientes/:id                     & DELETE & Eliminar cliente                   \\
  /api/clientes/:id/documentos          & GET    & Facturas/cotizaciones del cliente  \\
  /api/clientes/report                  & GET    & Reporte de clientes                \\
  \midrule
  \multicolumn{3}{l}{\textit{Servicios (protegidas)}} \\
  /api/servicios                        & GET    & Listar servicios                   \\
  /api/servicios                        & POST   & Crear servicio                     \\
  /api/servicios/:id                    & PUT    & Actualizar servicio                \\
  /api/servicios/:id                    & DELETE & Eliminar servicio                  \\
  /api/servicios/categorias             & GET    & Listar categorías de servicios     \\
  \midrule
  \multicolumn{3}{l}{\textit{Gastos (protegidas)}} \\
  /api/gastos                           & GET    & Listar gastos                      \\
  /api/gastos                           & POST   & Crear gasto                        \\
  /api/gastos/:id                       & PUT    & Actualizar gasto                   \\
  /api/gastos/:id                       & DELETE & Eliminar gasto                     \\
  /api/gastos/categorias                & GET    & Listar categorías de gastos        \\
  \midrule
  \multicolumn{3}{l}{\textit{Reportes y Finanzas (protegidas)}} \\
  /api/reportes/resumen                 & GET    & Resumen general de ventas          \\
  /api/reportes/por-fecha               & GET    & Reporte de ventas por período      \\
  /api/reportes/por-cliente             & GET    & Reporte de ventas por cliente      \\
  /api/reportes/por-servicio            & GET    & Reporte de ventas por servicio     \\
  /api/finanzas/resumen                 & GET    & Resumen financiero consolidado     \\
  \midrule
  \multicolumn{3}{l}{\textit{Configuración (protegida)}} \\
  /api/config                           & GET    & Obtener configuración de empresa   \\
  /api/config                           & PUT    & Actualizar configuración           \\
  \midrule
  \multicolumn{3}{l}{\textit{Correo Outlook OAuth2 (protegidas)}} \\
  /api/email/outlook/init               & POST   & Iniciar flujo de autorización      \\
  /api/email/outlook/poll               & POST   & Consultar estado de autorización   \\
  /api/email/outlook/status             & GET    & Verificar conexión activa          \\
  /api/email/outlook/disconnect         & DELETE & Desconectar cuenta Outlook         \\
  \midrule
  \multicolumn{3}{l}{\textit{Respuesta pública de cotizaciones (sin autenticación)}} \\
  /api/cotizacion-respuesta/:token      & GET    & Obtener cotización por token       \\
  /api/cotizacion-respuesta/:token      & POST   & Aprobar o rechazar cotización      \\
  \bottomrule
\end{longtable}

% ─────────────────────────────────────────────────────────────────────────────
\section{Seguridad y Control de Datos}

\begin{itemize}
  \item \textbf{Autenticación JWT:} Todas las rutas protegidas requieren un
        token Bearer válido. Los tokens se generan al iniciar sesión y tienen
        expiración configurable.
  \item \textbf{Contraseñas cifradas:} Las contraseñas de usuarios se almacenan
        con hash bcrypt; nunca en texto plano.
  \item \textbf{Tokens de aprobación:} Los tokens de aprobación de cotizaciones
        son aleatorios de 32 bytes (256 bits), únicos por cotización y con
        expiración de 7 días. Se invalidan tras ser usados.
  \item \textbf{Eliminación lógica (Soft Delete):} Los documentos anulados
        no se eliminan físicamente; se marca \texttt{status = `anulada'}.
  \item \textbf{Credenciales SMTP protegidas:} El endpoint público de
        cotizaciones excluye las credenciales SMTP de la respuesta.
  \item \textbf{Integridad referencial:} Una factura generada desde cotización
        almacena el ID y número de la cotización original para trazabilidad.
  \item \textbf{Consecutivo correlativo:} El sistema incrementa automáticamente
        el número de factura dentro del rango autorizado.
\end{itemize}

% ─────────────────────────────────────────────────────────────────────────────
\section{Generación de PDF e Impresión}

\begin{itemize}
  \item El sistema genera el PDF mediante la función de impresión del
        navegador web (\textit{Ctrl+P $\rightarrow$ Guardar como PDF}).
  \item El formato de página es \textbf{carta (Letter)} en orientación
        vertical, con márgenes configurados vía CSS \texttt{@page}.
  \item La hoja impresa incluye todos los campos obligatorios del SAR:
        CAI, RTN, rango autorizado, fecha límite de emisión, totales
        desglosados e ISV.
  \item Las cotizaciones muestran el RTN del emisor pero \textbf{no muestran}
        CAI ni información de rango autorizado, al ser documentos pre-venta.
  \item El cliente puede imprimir la cotización directamente desde el
        enlace de aprobación enviado por correo, visualizando el documento
        formal completo.
\end{itemize}

% ─────────────────────────────────────────────────────────────────────────────
\section{Requisitos para Registro y Autorización SAR}

\subsection{Documentos Legales de la Empresa}

\begin{enumerate}[label=\textbf{\arabic*.}]
  \item Escritura de constitución de la empresa (copia autenticada).
  \item RTN de la empresa (vigente).
  \item RTN del representante legal.
  \item Copia de identidad del representante legal.
  \item Poder de representación (si aplica).
  \item Constancia de inscripción en el SAR.
  \item Declaración de Impuesto Sobre Ventas (ISV) al día.
  \item Declaración del Impuesto Sobre la Renta (ISR) al día.
\end{enumerate}

\subsection{Información del Sistema de Facturación}

\begin{enumerate}[label=\textbf{\arabic*.}]
  \item \textbf{Tipo de sistema:} Electrónico (aplicación web).
  \item \textbf{Tecnología:} Node.js + React (JavaScript).
  \item \textbf{Formato de correlativo:} \texttt{000-001-01-XXXXXXXX}.
  \item \textbf{Campos incluidos:} CAI, RTN emisor, RTN adquirente,
        fecha, rango, venta gravada, exonerada, exenta, ISV 15\%,
        total, suma en letras.
  \item \textbf{Respaldo de datos:} Archivos locales \texttt{.db}
        (NeDB); posibilidad de exportación a CSV/Excel.
  \item \textbf{Impresión/PDF:} Generación vía navegador web en
        formato carta.
\end{enumerate}

\subsection{Información para Solicitud de CAI}

\begin{tabularx}{\linewidth}{>{\bfseries}l X}
  \toprule
  Dato Requerido           & Descripción                                  \\
  \midrule
  Razón Social             & DEVS MRS S. DE R. L.                         \\
  RTN                      & RTN de la empresa                            \\
  Número de Establecimientos & 1 (sede principal)                          \\
  Cantidad de Folios       & Cantidad de facturas a autorizar             \\
  Serie / Correlativo Desde& Número inicial del rango (ej.\ 00000001)     \\
  Serie / Correlativo Hasta& Número final del rango (ej.\ 00001000)       \\
  Tipo de Documento        & Factura Comercial                            \\
  Modalidad de Facturación & Sistema Computarizado / Electrónico          \\
  \bottomrule
\end{tabularx}

% ─────────────────────────────────────────────────────────────────────────────
\section{Glosario de Términos}

\begin{tabularx}{\linewidth}{>{\bfseries}l X}
  \toprule
  Término    & Definición                                                   \\
  \midrule
  SAR        & Servicio de Administración de Rentas de Honduras             \\
  RTN        & Registro Tributario Nacional                                 \\
  CAI        & Código de Autorización de Impresión                          \\
  ISV        & Impuesto Sobre Ventas (15\%)                                 \\
  ISR        & Impuesto Sobre la Renta                                      \\
  JWT        & JSON Web Token (token de autenticación firmado digitalmente) \\
  OAuth2     & Protocolo de autorización delegada (usado para Outlook)      \\
  Graph API  & API de Microsoft para acceso a servicios de Microsoft 365    \\
  SMTP       & Simple Mail Transfer Protocol (protocolo de envío de correo) \\
  SPA        & Single Page Application (aplicación web de página única)    \\
  API REST   & Interfaz de programación basada en HTTP con JSON             \\
  NeDB       & Node Embedded Database (base de datos embebida en Node.js)  \\
  HNL        & Lempira Hondureño (moneda nacional)                          \\
  \bottomrule
\end{tabularx}

% ─────────────────────────────────────────────────────────────────────────────
\vspace{1cm}
\begin{center}
  \rule{0.5\linewidth}{0.5pt}\\[0.3cm]
  {\small Documento generado por DEVS MRS S.\ de R.\ L.\ --- Mayo 2026\\
  Para uso en trámites de registro y autorización ante el SAR de Honduras.\\
  Versión 2.0 --- Incluye módulos de clientes, servicios, gastos, reportería,
  finanzas, envío de cotizaciones por correo y autenticación de usuarios.}
\end{center}

\end{document}
